\documentclass[11pt,a4paper]{article}
\usepackage[margin=1.05in]{geometry}
\usepackage{amsmath,amssymb,amsthm}
\usepackage[hidelinks]{hyperref}
\newtheorem{theorem}{Theorem}
\newtheorem{lemma}[theorem]{Lemma}
\newtheorem{proposition}[theorem]{Proposition}
\newcommand{\floor}[1]{\lfloor #1\rfloor}
\newcommand{\ceil}[1]{\lceil #1\rceil}
\hypersetup{pdftitle={A quadratic alphabet threshold for non-overlapping codes},pdfauthor={Peiwei Wu and Qian Peng}}
\title{A quadratic alphabet threshold for non-overlapping codes}
\author{Peiwei Wu$^{1}$ \qquad Qian Peng$^{2,*}$\\[0.5ex]
\small $^{1}$Shandong Medical and Pharmaceutical University\\
\small (formerly Binzhou Medical University)\\
\small $^{2}$The University of Melbourne\\
\small $^{*}$Corresponding author}
\date{}
\begin{document}
\maketitle
\begin{abstract}
Let $q_0(n)$ be the least integer such that every alphabet of size $q\ge q_0(n)$ admits a partition $I\sqcup J$ for which $I^{n-1}J$ is a maximum non-overlapping code of length $n$. Chee proved that $q_0(n)$ exists. We prove $q_0(n)=n^2/2+O(n)$ and the finite bound $q_0(n)\le4n(n-1)$. The upper bound compares a uniform loss in the real partition relaxation with the cost of rounding a single alphabet split. Moving one two-letter word between parts produces an infinite family of exceptions and disproves a conjecture of Stanovnik, Mo\v{s}kon and Mraz. We determine the maximum for every alphabet size at lengths five and six. Their thresholds are $4$ and $10$, respectively. At length six the exceptional alphabet sizes are $2,3,4,9$. On nine symbols the maximum is $33{,}872$, exceeding the best $I^5J$ code by 258.
\end{abstract}

\section{The alphabet threshold}

A code $C\subseteq\Sigma^n$ is \emph{non-overlapping} if a nonempty proper prefix of a word in $C$ never equals a suffix of another word in $C$ or of the same word. Let $S_n(q)$ denote the largest size of such a code when $|\Sigma|=q$. For any partition $\Sigma=I\sqcup J$ into nonempty parts, $I^{n-1}J$ is non-overlapping. The best size given by these Blackburn codes is
\[
 B_n(q)=\max_{1\le s<q}s^{n-1}(q-s).
\]
Blackburn conjectured that $S_n(q)=B_n(q)$ for each fixed $n\ge2$ and all sufficiently large $q$~\cite{Blackburn}. Chee proved this conjecture and asked for the least integer $q_0(n)\ge2$ such that equality holds for every $q\ge q_0(n)$~\cite[Theorem~I.1 and Section~IV]{Chee}.

\begin{theorem}\label{th:threshold}
For all $n\ge3$, $q_0(n)\le 4n(n-1)$. For every $n\ge12$,
\[
 q_0(n)>n\bigl(\floor{n/2}-2\bigr)+\ceil{n/2}.
\]
For $n\ge564$, $q_0(n)\le\floor{n^2/2+143n}+1$. In particular, $q_0(n)=n^2/2+O(n)$.
\end{theorem}

Stanovnik, Mo\v{s}kon and Mraz conjectured that some maximum partition code has $R_i=\varnothing$ for all $i>1+n/(q+1)$, in the notation of~\eqref{eq:levels}. For $q\ge n$ this forces every right part above the first level to be empty and gives $S_n(q)=B_n(q)$~\cite[Conjecture~1 and Remark~3]{SMM}. An example of length six already contradicts that prediction.

At lengths two, three and four the maximum equals $B_n(q)$ over every alphabet~\cite{Blackburn,SMM}. The next two lengths admit complete classifications.
\begin{theorem}\label{th:five}
$S_5(2)=2$, $S_5(3)=17$, and $S_5(q)=B_5(q)$ for all $q\ge4$. Thus $q_0(5)=4$.
\end{theorem}
\begin{theorem}\label{th:six}
For $q\ge2$, one has $S_6(q)=B_6(q)$ except at four alphabet sizes, whose maxima are given below.
\[
\begin{array}{c|rrrr}
q&2&3&4&9\\\hline
S_6(q)&3&41&251&33{,}872\\
B_6(q)&1&32&243&33{,}614
\end{array}
\]
In particular, $q_0(6)=10$.
\end{theorem}

Table~1 of~\cite{SMM} already gives the maxima for $2\le q\le6$ at both lengths, including the first three exceptional values in Theorem~\ref{th:six}. The accompanying data give an optimal ternary partition at length five with a single-word split at level two, up to reversal~\cite{SMMdata}. Theorems~\ref{th:five} and~\ref{th:six} extend these tables to every alphabet size.

\section{An exceptional alphabet for every large length}

We use the construction of Fimmel et al.~\cite{Fimmel} in the form analyzed by Stanovnik, Mo\v{s}kon and Mraz~\cite[Construction~6]{SMM}. Partition $W_1=\Sigma$ into $L_1,R_1$. For each $i\ge2$, form
\begin{equation}\label{eq:levels}
 W_i=\bigcup_{j=1}^{i-1}L_jR_{i-j},
 \qquad W_i=L_i\sqcup R_i\quad (i<n).
\end{equation}
The union over $j$ is disjoint. Every $W_n$ constructed in this way is non-overlapping, and every maximum code has such a representation~\cite[Theorems~3 and 6, Proposition~7]{SMM}. Reversing the words interchanges the two sides. Write $l_i,r_i,w_i$ for the respective cardinalities.

Set $l_1=a$, $r_1=b$ and move a single word of $W_2=L_1R_1$ into $L_2$. Put every other two-letter word and every subsequent level on the right. Thus $r_1=b$, $r_2=ab-1$, and the resulting code has size $r_n$, where
\[
 r_j=ar_{j-1}+r_{j-2}\qquad (3\le j\le n).
\]
We use this same family for all lengths in the lower bound.

\begin{proof}[Proof of the lower bound in Theorem~\ref{th:threshold}]
Let $k=\floor{n/2}-2$, $d=\ceil{n/2}$, $q=nk+d$, $b=k+1$ and $a=q-b$. Then $(n-2)b-a=1$ and $a>n-2$ for $n\ge12$. The continuous maximum of $f(s)=s^{n-1}(q-s)$ occurs at $s_*=q(n-1)/n$, with $s_*-a=1-d/n\in(0,1/2]$. In fact $a$ is an integer maximizer. The identity $f'(s)=ns^{n-2}(s_*-s)$ shows that for equal distances from $s_*$ the loss to the right is at least the loss to the left. Since $a+1$ is at least as far from the peak as $a$, $f(a+1)\le f(a)$, and unimodality handles the other integers.

Write $v_j=r_j/a^{j-1}$. Then $v_1=b$, $v_2=b-1/a$ and $v_j=v_{j-1}+v_{j-2}/a^2$. All terms are at least $b-1/a$. Summing the increments from $j=3$ to $n$ yields
\[
 v_n\ge b-\frac1a+\frac{(n-2)(b-1/a)}{a^2}
     = b+\frac{1}{a^2}-\frac{n-2}{a^3}>b .
\]
It follows that $r_n>ba^{n-1}=B_n(q)$, so $q_0(n)>q$. Dividing this choice of $q$ by $n^2$ gives $\liminf_{n\to\infty} q_0(n)/n^2\ge1/2$.
\end{proof}

\section{A uniform loss for moderate partitions}

Normalize each level cardinality in~\eqref{eq:levels} by dividing it by $q^i$ and retain the notation $l_i,r_i,w_i$ for the normalized quantities. Write $a=l_1$, $b=r_1$, $M=w_n$; thus $a+b=1$. Put
\[
 c=\frac{n-1}{n},\qquad m=\frac{1}{n-1},\qquad
 \delta=\frac{1}{2n(n-1)},\qquad \beta(x)=x^{n-1}(1-x).
\]
If $\max(a,b)\ge (n-2)/(n-1)$, Chee's skew-partition theorem gives $M\le B_n(q)/q^n$~\cite[Theorem~III.2]{Chee}. For the remaining partitions, both $a$ and $b$ lie in $[m,c-2\delta]$. We call these partitions moderate; the following estimate also holds at the endpoints.

Let $L(z)=\sum_{i=1}^{n-1} l_i z^i$, let $\rho>0$ satisfy $L(\rho)=1$, and put $t=1/\rho$. Chee's root identity~\cite[Lemma~II.5]{Chee} reads
\begin{equation}\label{eq:root}
 M=\beta(t)-D,\qquad
 D=\sum_{\substack{1\le i,j<n\\i+j>n}} l_i r_j\rho^{i+j-n}.
\end{equation}
For a nonempty code $M>0$, so $0<t<1$ and $\rho>1$. The identity is valid in either orientation of the partition, using the corresponding $L$ and $\rho$.

\begin{lemma}\label{lem:gap}
If $q\ge2n$ and a moderate partition attains $S_n(q)$, then
\[
 \beta(c)-M\ \ge\ \frac{1}{96n(n-1)^2}.
\]
\end{lemma}
\begin{proof}
First, choose $s=\ceil{qc}$. Then $s/q\ge c$ and $1-s/q\ge1/(2n)$. Since $c^{n-1}\ge1/3$, the Blackburn code of this split has normalized size at least $1/(6n)$, so $M\ge1/(6n)$. We use $(1+1/j)^j\le e<3$ for positive integers $j$.

Suppose the endpoint terms $ar_{n-1}+bl_{n-1}$ in $w_n$ account for at least $M/2$. Since $a+b=1$, one of $r_{n-1},l_{n-1}$ is at least $M/2$. Reverse the partition if needed to arrange $r_{n-1}\ge M/2$, then take the root in~\eqref{eq:root}. The terms with $j=n-1$ give
\[
 D\ge r_{n-1}\sum_{i=2}^{n-1} l_i\rho^{i-1}
       = r_{n-1}(t-a).
\]
If $t\ge c-\delta$, then $t-a\ge\delta$ and $\beta(c)-M\ge M\delta/2$. If $t<c-\delta$, integrate $\beta'(x)=nx^{n-2}(c-x)$ from $c-\delta$ to $c$. Since $(c-\delta)^{n-2}\ge((n-2)/(n-1))^{n-2}\ge1/3$, monotonicity on $[0,c]$ gives $\beta(c)-M\ge n\delta^2/6$.

Otherwise the terms $l_i r_j$ with $i+j=n$ and $i,j\ge2$ sum to at least $M/2$. For each such pair, the recurrence yields $w_{i+1}\ge m l_i$ and $w_{j+1}\ge m r_j$. If $l_{i+1}\ge ml_i/2$, select the tail term $l_{i+1}r_j$. Otherwise, if $r_{j+1}\ge mr_j/2$, select $l_i r_{j+1}$. In the remaining case $r_{i+1}\ge ml_i/2$ and $l_{j+1}\ge mr_j/2$, so select $l_{j+1}r_{i+1}$. In every case the selected term is at least $m^2 l_i r_j/4$. Each term of total degree $n+1$ can be selected at most twice, by decrementing its first or its second index. Each term of degree $n+2$ can be selected at most once, since the ordered pair $(j+1,i+1)$ determines $(i,j)$. The weights in $D$ are at least one, so $D\ge m^2 M/16$.

The three lower bounds $M\delta/2$, $n\delta^2/6$ and $m^2M/16$ are each at least $1/(96n(n-1)^2)$.
\end{proof}

\begin{proof}[Proof of the upper bound in Theorem~\ref{th:threshold}]
Assume $q\ge4n(n-1)$ and round $qc$ to a nearest integer $s$, which lies between 1 and $q-1$. Since $|\beta'(x)|\le n|x-c|$ on $[0,1]$,
\[
 \beta(c)-\beta(s/q)\le\frac{n}{8q^2}<
 \frac{1}{96n(n-1)^2}.
\]
The Blackburn code with split $s$ thus beats every putative moderate maximum, by Lemma~\ref{lem:gap}. A maximum must have a skew partition, for which Chee's theorem gives the asserted equality.
\end{proof}

\section{The linear error term}

\begin{proof}[Proof of the asymptotic formula in Theorem~\ref{th:threshold}]
The alphabet sizes in the lower bound differ from $n^2/2$ by $O(n)$. We prove that any exceptional size $q\ge n^2/2$ also differs from $n^2/2$ by at most $O(n)$, with absolute constants throughout. Choose a maximum partition for such an alphabet. Since $S_n(q)>B_n(q)$, Chee's skew-partition theorem makes it moderate. Put
\[
 \Delta=\beta(c)-M,\qquad P=c^{n-2},\qquad
 F=\sum_{i=2}^{n-2} l_i r_{n-i}.
\]
The root identity gives $0\le D\le\Delta$. Let $s$ be a nearest integer to $qc$. The rounding estimate of Section~3 gives $\Delta\le n/(8q^2)\le1/(2n^3)$.

For $i+j=n$, $i,j\ge2$, use $w_{i+1}\ge b l_i$ and $w_{j+1}\ge a r_j$ in the interior argument of Lemma~\ref{lem:gap}. Select $l_{i+1}r_j$ if $l_{i+1}\ge b l_i/2$, and otherwise select $l_i r_{j+1}$ if $r_{j+1}\ge a r_j/2$. In the remaining case $l_{j+1}r_{i+1}\ge ab l_i r_j/4$. Each choice is at least $ab l_i r_j/4$, with multiplicity at most two. Hence $D\ge abF/8$. Moderateness gives $ab\ge m(1-m)\ge1/[2(n-1)]$, so $F\le16(n-1)\Delta=O(n^{-2})$.

The endpoint terms sum to $M-F=a r_{n-1}+b l_{n-1}$. Reverse the partition if necessary so that $r_{n-1}\ge M-F$, and use the root $\rho$ of this orientation. Reversal preserves $M$ and $F$. The root equation gives $t\ge a$, and the terms with $j=n-1$ in~\eqref{eq:root} give $D\ge(M-F)(t-a)$.

Write $u=n^2(c-t)$. For $n\ge4$ we have $|u|\le3$. Indeed, if $t<c-1/n$, monotonicity and integration of $\beta'(x)=n x^{n-2}(c-x)$ on $[c-1/n,c]$ give
\[
 \Delta\ge\beta(c)-\beta(t)\ge\frac{(1-2/n)^{n-2}}{2n}\ge\frac1{18n},
\]
since $(1+2/(n-2))^{n-2}\le e^2<9$. This contradicts $\Delta\le1/(2n^3)$. For $t\ge c-1/n$, every $x$ between $t$ and $c$ satisfies $x^{n-2}\ge1/9$. Integration on that interval gives $n(c-t)^2/18\le\Delta$, proving the claim.

Both $t$ and $s/q$ lie within $3/n^2$ of $c$. For every $x$ between $c$ and either point, the mean value theorem on $[0,1]$ gives
\[
 |x^{n-2}-P|\le(n-2)|x-c|\le3/n .
\]
Since $P\ge c^{n-1}\ge1/3$ and $|u|\le3$, integration gives
\begin{align}
 \frac{n^3}{P}\bigl(\beta(c)-\beta(t)\bigr)
 &\ge\left(1-\frac9n\right)\frac{u^2}{2},\label{eq:sharploss}\\
 \frac{n^3\Delta}{P}
 &\le \frac{n^4}{8q^2}\left(1+\frac3n\right).\label{eq:sharpgrid}
\end{align}
In the grid estimate the integration interval has length at most $1/(2q)\le1/n^2$, giving $|x^{n-2}-P|\le1/n$ there.

Set $w=n(M-F)/P$. The identities $1-c=1/n$ and $nM/P=c-n\Delta/P$ give
\[
 1-w=\frac1n+\frac{n(\Delta+F)}P\le\frac{25}{n},
\]
since $\Delta+F\le16n\Delta$ and $\Delta\le1/(2n^3)$. In particular $0\le w\le1$. Moderateness gives $a\le c-1/[n(n-1)]$, so $n^2(t-a)\ge n/(n-1)-u\ge1-u$. The endpoint estimate now gives $n^3D/P\ge w(1-u)\ge 1-u-100/n$, since $|1-u|\le4$. Adding~\eqref{eq:sharploss} yields
\[
 \frac{n^3\Delta}{P}\ge\frac{u^2}{2}+1-u-\frac{141}{n}
 =\frac{1+(u-1)^2}{2}-\frac{141}{n}.
\]
If $n\ge564$, comparison with~\eqref{eq:sharpgrid} gives
\[
 \frac{2q}{n^2}\le\sqrt{1+\frac{285}{n-282}}
             \le 1+\frac{285}{2(n-282)}\le 1+\frac{285}{n}.
\]
In particular, $q\le n^2/2+143n$. All larger alphabets admit a maximum Blackburn code, proving the finite bound and the asymptotic formula.
\end{proof}

\section{Exact optimization at length five}

\begin{proof}[Proof of Theorem~\ref{th:five}]
Return to integer counts and use reversal to assume $a=l_1\ge b=r_1>0$, $q=a+b$. Write $x=l_2$, $y=r_2=ab-x$, and $T=ay+bx=w_3$. With $u=l_3$ and $s=l_4$ the last two levels obey
\[
 w_4=a(T-u)+xy+bu,\qquad
 w_5=a(w_4-s)+x(T-u)+uy+bs.
\]
At fixed $x,u$ the coefficient of $s$ is $b-a\le0$, so $s=0$ maximizes $w_5$. The remaining expression is affine in $u$. Its two endpoint values, at $u=0,T$ respectively, are
\begin{align*}
 F_0(x)&=a^4b+a^2(3b-a)x-(2a-b)x^2,\\
 F_1(x)&=2a^3b^2-2ab(a-b)x-bx^2 .
\end{align*}
The maximum for the split is therefore $\max_{0\le x\le ab}\max(F_0(x),F_1(x))$, with $x$ integral.

When $a\ge3b$, both polynomials are at most $a^4b$. When $b\le a\le3b$, both are at most $81q^5/1024$. For $F_1$ this follows by bounding it by $2a^3b^2$, whose maximum for real $a+b=q$ is $216q^5/3125$. For $F_0$, complete the square to obtain
\[
 F_0(x)\le\frac{a^4((a+b)^2+4b^2)}{4(2a-b)}.
\]
To compare this with $81q^5/1024$, multiply through by the positive denominators. The required nonnegative difference factors as
\[
\begin{split}
 &81(2a-b)(a+b)^5-256a^4((a+b)^2+4b^2) \\
 &\quad=(3b-a)
 \bigl(94a^5+65a^4b+260a^3b^2-30a^2b^3-90ab^4-27b^5\bigr).
\end{split}
\]
The three subtracted terms have sum at most $147a^3b^2$.

If $[3q/4,4q/5]$ contains an integer $s$, then $s^4(q-s)\ge81q^5/1024$ by monotonicity up to $4q/5$. This condition holds whenever $q=4h+5j$ with $h,j\ge0$, by taking $s=3h+4j$. It covers $q=4,5,8,9,10$ and every $q\ge12$ (use $12,13,14,15$ and add fours). For $q=6,11$ use $s=5,9$ respectively; direct substitution gives the same lower bound. The remaining small splits are displayed below. Each $F_0$ maximum is obtained from the integer or integers nearest its vertex in $[0,ab]$, and $F_1$ is maximized at zero.
\[
\begin{array}{c|c|r|r}
q&a&\max F_0&\max F_1\\\hline
2&1&2&2\\
3&2&17&16\\
7&4&1088&1152\\
7&5&1268&1000\\
7&6&1296&432
\end{array}
\]
For $q=7$ the maximum is $6^4=B_5(7)$. Every endpoint partition used in the calculation is feasible in~\eqref{eq:levels}, which gives equality in the two exceptional cases as well.
\end{proof}

\section{Exact optimization at length six}
\begin{proof}[Proof of Theorem~\ref{th:six}]
We first realize the exceptional sizes. For $q=2$ use the words $001011$, $001101$ and $001111$. Their one-letter prefixes and suffixes differ. At lengths two through five, each prefix starts with $00$, which occurs only at the start of each word. For $q=3,4$ use the single marked pair of Section~2 with $a=q-1$ and $b=1$. Then $r_2=a-1$ and $r_j=ar_{j-1}+r_{j-2}$, giving $r_6=41,251$, respectively.

For $q=9$, let $I=\{0,\ldots,6\}$ and $J=\{7,8\}$. Take all words starting in $I$ and ending in $J$ such that every interior letter from $J$ is a 7 immediately preceded by 0 or 1. Exclude the two endings 07 and 17. Any proper prefix ending in $J$ ends with one of these pairs, so it cannot equal a suffix. The first two counts are $r_1=2$ and $r_2=12$. Removing either one initial letter from $I$ or one of the two marked pairs gives $r_j=7r_{j-1}+2r_{j-2}$. Thus $r_6=33{,}872$. With just the pair 07 the recurrence in Section~2 gives $33{,}851$, which already exceeds $B_6(9)$ and refutes Conjecture~1 of~\cite{SMM}.

For the upper bound it suffices by Theorem~\ref{th:threshold} to treat $2\le q<120$. Reverse the partition if necessary to ensure $a=l_1\ge b=r_1>0$. Put $x=l_2$, $y=ab-x$, $T=ay+bx=w_3$ and $u=l_3$. Then
\[
 w_4=a(T-u)+xy+bu.
\]
At level five we may place all words on the right, since the coefficient of $l_5$ in $w_6$ is $b-a\le0$. For a fixed $u$, a word assigned to $R_4$ has coefficient $a^2+x$ in $w_6$, whereas one assigned to $L_4$ has coefficient $ab+y$. Hence with $D=\max(a^2+x,ab+y)$ the largest objective for these three initial splits is
\begin{align*}
 D w_4+a\bigl(x(T-u)+uy\bigr)+u(T-u)
    &= A+Ku-u^2,\\
 A &= D(aT+xy)+axT,\\
 K &= D(b-a)+a(y-x)+T.
\end{align*}
For nonnegative integral $u$ this is at most $A+\floor{\max(K,0)^2/4}$. The following integer calculation evaluates that upper bound at every remaining $q,a,x$ and records just the strict excesses over $B_6(q)$.
\begin{verbatim}
exceptions = {}
for q in range(2, 120):
    B = max(s**5 * (q-s) for s in range(1, q))
    U = 0
    for a in range((q+1)//2, q):
        b = q-a
        for x in range(a*b+1):
            y = a*b-x
            T = a*y+b*x
            D = max(a*a+x, a*b+y)
            A = D*(a*T+x*y)+a*x*T
            K = D*(b-a)+a*(y-x)+T
            U = max(U, A+max(K,0)**2//4)
    if U > B:
        exceptions[q] = U
\end{verbatim}
The four entries are $(2,3),(3,41),(4,251),(9,33872)$. The constructions above attain all four bounds. For every other $q$ a Blackburn split attains the upper bound.
\end{proof}

\paragraph{Formal verification.}
The accompanying Lean~4 development defines a code as a finite set of words, constructs the Blackburn codes and the exceptional family of Section~2, and verifies the bounds and limiting ratio for the least threshold $q_0(n)$. It also certifies the maximum at lengths five and six for every alphabet size. For the finite calculation at length six, it bounds the quadratic envelope on intervals of $x$, subdividing any interval until its bound is at most the claimed maximum.

For the upper bound the formal proof constructs integer layer counts that dominate each actual code $C$. Begin with each codeword split into single letters. Label a block left if it is a prefix of a word of $C$ and right otherwise. While at least two blocks remain, the first is left and the last is right by non-overlap. Merge an adjacent left--right pair and repeat. This parses every codeword. Counting possible pairs bounds the number of parsed words of length $i$ by the convolution of the shorter left and right counts. Keep the left counts and increase the right counts successively until equality holds in the recurrence. Induction then bounds $|C|$ by the final convolution. This argument permits multiple parses of a word. Lean applies the verified root identity, skew estimate and moderate-partition bounds to the resulting integer profile.

\paragraph{Use of generative AI.}
OpenAI Codex was used to develop the arguments, write the Lean formalization and draft the manuscript.

\begin{thebibliography}{9}
\raggedright
\bibitem{Blackburn}
S.~R. Blackburn, Non-overlapping codes, \emph{IEEE Trans. Inform. Theory} \textbf{61} (2015), 4890--4894. \href{https://doi.org/10.1109/TIT.2015.2456634}{doi:10.1109/TIT.2015.2456634}.
\bibitem{Chee}
Y.~M. Chee, A proof of Blackburn's conjecture on non-overlapping codes, 2026. \href{https://arxiv.org/abs/2611.28776v2}{arXiv:2611.28776v2}.
\bibitem{Fimmel}
E.~Fimmel, C.~J. Michel, F.~Pirot, J.-S. Sereni and L.~Str\"ungmann, \href{https://ajc.maths.uq.edu.au/pdf/86/ajc_v86_p233.pdf}{Diletter and triletter comma-free codes over finite alphabets}, \emph{Australas. J. Combin.} \textbf{86} (2023), 233--270.
\bibitem{SMM}
L.~Stanovnik, M.~Mo\v{s}kon and M.~Mraz, In search of maximum non-overlapping codes, \emph{Des. Codes Cryptogr.} \textbf{92} (2024), 1299--1326. \href{https://doi.org/10.1007/s10623-023-01344-z}{doi:10.1007/s10623-023-01344-z}.
\bibitem{SMMdata}
L.~Stanovnik, M.~Mo\v{s}kon and M.~Mraz, \href{https://github.com/magdevska/nono-codes}{Computational data for \emph{In search of maximum non-overlapping codes}}, dataset accompanying~\cite{SMM}.
\end{thebibliography}
\end{document}
